\documentclass[a4paper]{article}

\usepackage[spanish]{babel}
\usepackage{graphicx}
\usepackage{amsmath, amssymb}
\usepackage[margin=2cm]{geometry}
\usepackage{fancyhdr}
\usepackage{enumerate}
\usepackage[shortlabels]{enumitem}
\usepackage{parskip}
\usepackage[most]{tcolorbox}
\usepackage[hidelinks]{hyperref}
\usepackage{float}
\usepackage{booktabs}
\usepackage{mathrsfs}
\usepackage{tikz}
\usepackage{xcolor}

% cabecera
\pagestyle{fancy}
\fancyhead[l]{Astroestadística}
\fancyhead[c]{Quiz 3}
\fancyhead[r]{\today}
\fancyfoot[c]{\thepage}
\renewcommand{\headrulewidth}{0.2pt}

\begin{document}

\section{Quiz de Probabilidad y Estadística}

Seleccione la opción correcta en cada pregunta.

\begin{enumerate}

\item Se obtiene un intervalo de confianza del 95\% para la media poblacional de la magnitud absoluta de cierto tipo de estrellas: $[-3.42,-2.85]$. ¿Qué indica este nivel de confianza?
\begin{enumerate}[a)]
    \item Existe una probabilidad de 0.95 de que la media esté entre $-3.42$ y $-2.85$.
    \item Hay un 5\% de probabilidad de que una estrella tenga una magnitud absoluta fuera del intervalo.
    \item El 95\% de las estrellas de la población tiene magnitud absoluta entre $-3.42$ y $-2.85$.
    \item \textcolor{red}{El 95\% de los intervalos construidos de esta manera contiene la media real.}
\end{enumerate}
\smallskip
\textit{Justificación:} En el enfoque frecuentista, el 95\% describe la proporción de intervalos que contendrían la media real en muestreos repetidos.

\item Si un intervalo de confianza para la media de una población es $[50,60]$, ¿cuál es el error estándar si el nivel de confianza es del 95\% y la distribución es simétrica?
\begin{enumerate}[a)]
    \item \textcolor{red}{$2.55$}
    \item $5$
    \item No se puede determinar con esta información.
    \item $10$
\end{enumerate}
\smallskip
\textit{Justificación:} El margen de error es $5$; usando el valor crítico normal aproximado de $1.96$ para un 95\%, el error estándar es $5/1.96 \approx 2.55$.

\item Se desea estimar la media de una variable astronómica mediante un intervalo de confianza. Manteniendo fijo el nivel de confianza y la variabilidad poblacional, ¿qué sucede aproximadamente con el ancho del intervalo al cuadruplicar el tamaño de muestra?
\begin{enumerate}[a)]
    \item \textcolor{red}{Se reduce a la mitad.}
    \item Se cuadruplica.
    \item Se reduce a la cuarta parte.
    \item No cambia.
\end{enumerate}
\smallskip
\textit{Justificación:} El error estándar es proporcional a $1/\sqrt{n}$; al cuadruplicar $n$, el ancho del intervalo se reduce a la mitad.

\item Un investigador calcula un intervalo de confianza para la media de una población y cambia el nivel de confianza del 90\% al 99\%, manteniendo todo lo demás constante. ¿Qué ocurrirá con el intervalo?
\begin{enumerate}[a)]
    \item Tendrá el mismo ancho, pero cambiará su centro.
    \item Será más estrecho porque disminuye el nivel de significancia.
    \item \textcolor{red}{Será más ancho porque se requiere mayor probabilidad de cobertura.}
    \item Será más estrecho porque aumenta la precisión.
\end{enumerate}
\smallskip
\textit{Justificación:} Para aumentar la confianza se necesita un mayor valor crítico y, por tanto, un intervalo más ancho.

\item Se dispone de una muestra de 12 estrellas y se desea estimar la media de una población aproximadamente normal. La desviación estándar poblacional es desconocida. ¿Cuál distribución es apropiada para construir un intervalo de confianza para $\mu$?
\begin{enumerate}[a)]
    \item \textcolor{red}{$t$ de Student con 11 grados de libertad.}
    \item $t$ de Student con 12 grados de libertad.
    \item Chi-cuadrado con 11 grados de libertad.
    \item Normal estándar, independientemente del tamaño de muestra.
\end{enumerate}
\smallskip
\textit{Justificación:} Con desviación estándar poblacional desconocida y una muestra normal de tamaño $n=12$, se usa la distribución $t$ con $n-1=11$ grados de libertad.

\end{enumerate}
\bibliographystyle{plainurl}
\bibliography{bibliografia}
\end{document}